\chapter{Literature Survey}
This chapter reviews 38 published works on the techniques FL-IDS combines: federated learning and federated intrusion detection, robust aggregation against poisoned updates, two-stage classifier-plus-anomaly-detector designs, coupling of learning-based detection with SDN enforcement, and the datasets and evaluation practice of the field. It then identifies the gaps this project addresses.

\section{Overview}
Federated learning (FL) was introduced by McMahan \emph{et al.}~\cite{mcmahan2017fedavg} as a way to train a shared model across many clients by averaging locally trained weights (FedAvg), so raw data stays on the clients. Kairouz \emph{et al.}~\cite{kairouz2021} later collected the open problems of the area: non-IID data, communication cost, privacy and robustness to malicious participants. Intrusion detection was one of the first security applications of FL, and the area has grown to the point where Hernandez-Ramos \emph{et al.}~\cite{hernandez2025fliids} published a systematic review of FL-based IDS in \emph{ACM Computing Surveys}. Their review covers the models, datasets, aggregation functions and libraries used in this literature. It shows that a federated autoencoder or a federated deep classifier is a well-established idea in itself, and that the open problems lie in non-IID data, poisoning, and real deployment.

FedAvg is fragile when some clients are malicious. Blanchard \emph{et al.}~\cite{blanchard2017} showed that a single Byzantine client can steer the average arbitrarily, and Yin \emph{et al.}~\cite{yin2018trimmed} analysed coordinate-wise median and trimmed-mean aggregation and showed that they tolerate a bounded fraction of Byzantine clients, but break down once attackers reach the trimmed fraction, or 50\% for the median. Later work moved to the \emph{direction} of updates: comparing each update with a reference using cosine similarity.

The works below are grouped by theme. Section~\ref{sec:litsummary} then compares the closest systems with FL-IDS on the four capabilities that FL-IDS combines.

\section{Detailed Survey}

\subsection{Federated learning and federated intrusion detection}

\textbf{Hsu \emph{et al.}, 2019}~\cite{hsu2019noniid}. This paper proposes drawing each client's class proportions from a Dirichlet distribution as a controllable way to simulate non-identical client data in FL experiments. FL-IDS uses the same device to split Edge-IIoTset across its clients (Section~\ref{sec:datadesign}).

\textbf{Beutel \emph{et al.}, Flower, 2020}~\cite{beutel2020}. Flower is a framework for FL research that lets clients run as separate processes and lets the server strategy be replaced. FL-IDS builds its federated training on Flower, so that its clients are real separate processes and its trust-filtered aggregation is a custom strategy.

\textbf{Mothukuri \emph{et al.}, 2022}~\cite{mothukuri2022fl}. This work proposes federated anomaly detection for IoT networks using gated recurrent unit (GRU) models trained on decentralised device data, with an ensembler that combines updates from several sources. It shows that federated training can match centralised accuracy while keeping data on devices. It uses no supervised first stage, no defence against poisoned updates, and no network-level enforcement.

\textbf{Li \emph{et al.}, DeepFed, 2021}~\cite{li2021deepfed}. DeepFed builds an intrusion detection model for industrial cyber-physical systems from a convolutional neural network and a GRU, trained with FL across several industrial sites. Model parameters are protected with a Paillier-cryptosystem-based protocol. DeepFed protects updates in transit but does not filter malicious ones, and it has no zero-day backstop or SDN mitigation.

\textbf{Rey \emph{et al.}, 2022}~\cite{rey2022}. This paper applies FL to malware detection on IoT devices using the N-BaIoT dataset, with both a supervised multi-layer perceptron and an unsupervised autoencoder, and reports federated results close to centralised ones. It also shows that plain averaging is highly vulnerable to model poisoning, even with a single adversary, and that alternative aggregation functions improve resilience. This is direct evidence for the choice of a robust aggregator in FL-IDS. The work has no cascade and no network-level response.

\textbf{Nguyen and Beuran, FedMSE, 2025}~\cite{nguyen2025fedmse}. FedMSE combines a shrink autoencoder with a centroid one-class classifier in a semi-supervised federated setup and weights aggregation by the clients' reconstruction error. It is evaluated on N-BaIoT with Dirichlet non-IID splits and is reported to keep its performance when only half the gateways take part in a round. Like FL-IDS it trains an autoencoder federatedly on normal traffic, but it has no supervised first stage and no defence against malicious updates.

\textbf{Shahin \emph{et al.}, 2025}~\cite{shahin2025}. This paper proposes a two-stage hybrid FL framework that combines generative and predictive models for anomaly detection and classification, validated on N-BaIoT traffic from nine IoT devices. It is one of the few federated IDS works with a two-stage structure, but it evaluates no poisoning defence and no SDN enforcement.

\textbf{Sazid and Abbas, 2026}~\cite{sazid2026}. This preprint proposes a trust-aware federated hybrid IDS for edge computing. The server weights client models by a trust measure so that a compromised node submitting corrupted weights has less influence, and the authors evaluate it on public intrusion datasets. It shows the same concern as FL-IDS with poisoned updates, but it is a preprint and describes no SDN enforcement.

\subsection{Poisoning attacks and Byzantine-robust aggregation}

\textbf{Blanchard \emph{et al.}, Krum, 2017}~\cite{blanchard2017}. This paper shows that no linear aggregation rule, including averaging, tolerates even one Byzantine worker, and proposes Krum, which selects the single update closest to its neighbours. It is the starting point of distance-based robust aggregation.

\textbf{Yin \emph{et al.}, 2018}~\cite{yin2018trimmed}. Coordinate-wise median and trimmed mean are shown to reach optimal statistical rates under a bounded fraction of Byzantine clients, but they fail once attackers reach the trimmed fraction or 50\% for the median. FL-IDS applies a 15\% trimmed mean to the updates that survive its trust filter.

\textbf{Fang \emph{et al.}, 2020}~\cite{fang2020}. This paper designs local model poisoning attacks that deliberately push the global model away from its honest direction and applies them against Krum, trimmed mean and median aggregation, showing that these defences can be defeated by an informed attacker. It motivates testing FL-IDS against more than one attacker, including a stealthy one that is hard to separate from honest clients.

\textbf{Fung \emph{et al.}, 2020}~\cite{fung2020}. The authors study sybil attacks, in which one adversary controls many clients, show that existing Byzantine-tolerant schemes fall short against several of them, and propose FoolsGold, which flags sybils because their updates are unusually similar to each other. This is a reminder that update similarity can be used in both directions: honest clients are similar to the consensus, while colluding attackers may be too similar to each other.

\textbf{Cao \emph{et al.}, FLTrust, 2021}~\cite{cao2021fltrust}. FLTrust gives the server a small clean ``root'' dataset and computes its own reference update. Each client update gets a trust score from its cosine similarity to that reference, with negative similarity clipped to zero, and its magnitude is normalised to the reference. FL-IDS borrows the rule that an update pointing \emph{away} from the reference (cosine $<0$) is excluded. It uses the coordinate-wise median of the round's updates as the reference instead of a root dataset.

\textbf{Li \emph{et al.}, 2021}~\cite{li2021spatial}. This paper detects poisoned updates by clustering their geometry in parameter space and adapts the learning rate to the clients' behaviour over time with momentum-based speculation. The use of a client's history, not only the current round, is the idea behind the exponential moving average (EMA) trust score in FL-IDS.

\textbf{Xu \emph{et al.}, SignGuard, 2022}~\cite{xu2022signguard}. SignGuard shows that median- and distance-based defences can be circumvented by carefully crafted poisoning, and that the element-wise sign of an update together with magnitude and similarity statistics is effective for filtering malicious gradients. This supports pairing cosine similarity with a norm check, which FL-IDS does with norm clipping, against the sign-flipping attacker used in its evaluation.

\textbf{Li \emph{et al.}, 2024}~\cite{li2024experimental}. This experimental study evaluates eight aggregation schemes under five Byzantine attacks and proposes ClippedClustering. Its main finding is that with non-IID data some schemes fail even when no client is malicious, so the robustness of every scheme is limited. This is the same concern as research gap~3 below.

\textbf{Zhu \emph{et al.}, 2024}~\cite{zhu2024cosine}. The authors propose an angle-based aggregation, CKADA, argue that distance-based defences miss attackers who tamper with the direction of updates while angle-based ones lose accuracy on large models, and report better results than both families. It is the closest precedent for filtering on cosine similarity, which FL-IDS combines with a median-absolute-deviation threshold, norm clipping and EMA trust.

\textbf{Thein \emph{et al.}, 2024}~\cite{thein2024}. This work proposes a personalised federated IDS with a logit-adjustment loss for imbalanced clients and studies label-flipping and model-update poisoning, showing that a single malicious client can degrade the global model. It confirms that poisoning is a realistic threat for a federated IDS, but it has no cascade and no SDN part.

\textbf{Luong \emph{et al.}, Fed-LSAE, 2024}~\cite{luong2024}. Fed-LSAE compares the latent-space representation of each client's model, taken from the penultimate layer and an autoencoder, with the global one using centred kernel alignment, and excludes clients that look malicious. It is evaluated on CIC-ToN-IoT and N-BaIoT. It defends federated threat detection against poisoning, but it works on latent representations of the classifier, while FL-IDS filters weight updates of its autoencoder.

\subsection{Two-stage detection, autoencoders and boosting}

\textbf{Zhang \emph{et al.}, 2023}~\cite{zhang2023lgbmae}. This paper proposes a two-stage intrusion detector: LightGBM with focal loss makes a first decision, and the reconstruction residual of an autoencoder makes a second decision on samples the first stage called normal. It reports above 90\% accuracy, recall and F1 on NSL-KDD and UNSW-NB15. This is the direct precedent for the boosting-first cascade used in FL-IDS. FL-IDS differs in training the autoencoder federatedly and protecting it with a trust filter.

\textbf{Mirsky \emph{et al.}, Kitsune, 2018}~\cite{mirsky2018}. Kitsune is an online, unsupervised network IDS that uses an ensemble of small autoencoders and is light enough to run on a modest device. It shows that reconstruction error on normal traffic can detect attacks that were never labelled, which is the role of the autoencoder stage in FL-IDS.

\textbf{Ke \emph{et al.}, LightGBM, 2017}~\cite{ke2017lightgbm}. LightGBM is a gradient-boosted decision-tree library that uses gradient-based one-side sampling and exclusive feature bundling to train quickly on large tabular data. Tree ensembles are a strong baseline for tabular network-flow features, which FL-IDS confirms: on the same 500,000-row sample, LightGBM reached 0.887 macro F1 against 0.622 for a comparably sized neural network.

\textbf{Chen and Guestrin, XGBoost, 2016}~\cite{chen2016xgboost}. XGBoost is a scalable tree-boosting system with regularised objectives and sparsity-aware split finding. It is the other widely used boosting library. In the project's comparison of about 50 classifiers, XGBoost reached 0.983 macro F1, within 0.002 of the best model (random forest, 0.985); FL-IDS uses LightGBM in the cascade.

\subsection{Software-defined networking and enforcement}

\textbf{McKeown \emph{et al.}, OpenFlow, 2008}~\cite{mckeown2008}. OpenFlow lets a controller install match-action rules in the flow tables of switches, which separates the control plane from the data plane. A detection verdict can therefore become a drop rule at the first switch. FL-IDS uses OpenFlow~1.3 through its controller.

\textbf{Lantz \emph{et al.}, Mininet, 2010}~\cite{lantz2010}. Mininet emulates a network of hosts and switches on a single machine, so SDN experiments can run on real network software without hardware. It is the emulator used for the live-replay lab in FL-IDS, with Open vSwitch as the switch.

\textbf{El Houda \emph{et al.}, MiTFed, 2023}~\cite{elhouda2023mitfed}. MiTFed lets several network domains collaboratively train an attack detector with FL, uses SDN to enforce mitigation, and uses a blockchain to make the collaboration verifiable. It is the closest precedent for the ``FL decides, SDN enforces'' structure of FL-IDS. MiTFed does not use a two-stage cascade with an anomaly backstop.

\textbf{Pillai and Mathivanan, FL-NDR, 2026}~\cite{pillai2026}. FL-NDR treats the OpenFlow switches of an SDN as a sensing layer for federated detection, adds federated unlearning that rolls the global model back after a poisoned client is found, and installs OpenFlow rules at the border switches of the compromised domain. It is evaluated on an SDN DDoS dataset in simulation. It is close to FL-IDS in coupling FL, poisoning and SDN, but it is not a two-stage cascade and it focuses on DDoS.

\subsection{Datasets and evaluation practice}

\textbf{Ferrag \emph{et al.}, Edge-IIoTset, 2022}~\cite{ferrag2022edge}. This paper releases the dataset used in this project. It is collected on a seven-layer IoT/IIoT testbed with more than ten device types and covers 14 attacks in five groups (DoS/DDoS, information gathering, man-in-the-middle, injection and malware). The authors explicitly position it for both centralised and federated learning and report baseline results for both. Chapter~\ref{chap:method} shows that the released CSV has formatting artefacts which inflate results trained on it directly.

\textbf{Galal, 2026}~\cite{galal2026provenance}. This preprint shows that in Edge-IIoTset the ``field does not apply'' value is serialised as \texttt{0} or \texttt{0.0} depending on the source capture, so it encodes the label, and builds a leakage-free benchmark for precision-agriculture IDS using random stratified and leave-one-device-out splits. FL-IDS found and fixed the same artefact independently and later, and credits it to this work. Galal does not report the column-shifted MITM rows, and uses no flow features, connection-grouped split, federated training or SDN mitigation.

\textbf{Sharafaldin \emph{et al.}, CIC-IDS2017, 2018}~\cite{sharafaldin2018cicids}. This work introduced the CIC-IDS2017 dataset, whose features are computed per flow by CICFlowMeter (packet counts, rates, durations, payload statistics). Flow features of this kind are standard in network IDS. FL-IDS adopts the same idea, including CICFlowMeter's 120-second active timeout for splitting long sessions into flow records.

\textbf{Engelen \emph{et al.}, 2021}~\cite{engelen2021}. The authors re-examine CIC-IDS2017 and find problems in traffic generation, flow construction, feature extraction and labelling, with more than 20\% of the original traces reconstructed or relabelled, and show that such data-collection flaws can strongly affect model evaluation. This is the same kind of problem that FL-IDS found and repaired in Edge-IIoTset.

\textbf{Vaccari \emph{et al.}, MQTTset, 2020}~\cite{vaccari2020}. MQTTset is a dataset of MQTT traffic with several attack types, built for machine-learning IDS on IoT messaging. The project's earlier comparison of classical and deep-learning models, reported in the accompanying paper, used this dataset.

\textbf{Firouzi \emph{et al.}, DataSense, 2025}~\cite{firouzi2025}. DataSense provides synchronised sensor and network data from an IIoT testbed, with 50 attacks in seven categories, and a benchmark of machine- and deep-learning detectors with multi-objective feature selection. The same earlier comparison also used DataSense.

\textbf{Sommer and Paxson, 2010}~\cite{sommer2010}. This paper argues that machine learning is harder to apply to network intrusion detection than to most other domains, because of the open-ended range of normal and attack traffic, the high cost of false alarms and the difficulty of interpreting results. It explains why FL-IDS reports false alarms on normal traffic next to detection rates and measures unseen attack types explicitly.

\textbf{Arp \emph{et al.}, 2022}~\cite{arp2022}. The authors identify ten recurring pitfalls in machine learning for security, such as sampling bias, label inaccuracy, data snooping and spurious correlations, and show them in published work. Several of these apply to a random packet-level split of Edge-IIoTset, which is why FL-IDS splits by connection and fixes the label leak.

\section{Summary of Closest Systems}\label{sec:litsummary}
Table~\ref{tab:litcompare} compares the closest federated intrusion detection systems with FL-IDS on four capabilities: federated training, a two-stage cascade, a defence against poisoned updates, and enforcement through SDN. The table is based on the abstracts and main contributions of each work.

\begin{table}[H]
\centering
\renewcommand{\arraystretch}{1.25}
\small
\begin{tabular}{|p{4.2cm}|c|c|c|c|}
\hline
\textbf{System} & \textbf{FL} & \textbf{Two-stage} & \textbf{Poisoning defence} & \textbf{SDN} \\ \hline
Mothukuri \emph{et al.}~\cite{mothukuri2022fl} & \checkmark & -- & -- & -- \\ \hline
DeepFed~\cite{li2021deepfed} & \checkmark & -- & -- & -- \\ \hline
Rey \emph{et al.}~\cite{rey2022} & \checkmark & -- & evaluated & -- \\ \hline
FedMSE~\cite{nguyen2025fedmse} & \checkmark & -- & -- & -- \\ \hline
Shahin \emph{et al.}~\cite{shahin2025} & \checkmark & \checkmark & -- & -- \\ \hline
Sazid and Abbas~\cite{sazid2026} & \checkmark & -- & \checkmark & -- \\ \hline
Thein \emph{et al.}~\cite{thein2024} & \checkmark & -- & \checkmark & -- \\ \hline
Fed-LSAE~\cite{luong2024} & \checkmark & -- & \checkmark & -- \\ \hline
Zhang \emph{et al.}~\cite{zhang2023lgbmae} & -- & \checkmark & -- & -- \\ \hline
MiTFed~\cite{elhouda2023mitfed} & \checkmark & -- & -- & \checkmark \\ \hline
FL-NDR~\cite{pillai2026} & \checkmark & -- & \checkmark & \checkmark \\ \hline
\textbf{FL-IDS (this work)} & \checkmark & \checkmark & \checkmark & \checkmark \\ \hline
\end{tabular}
\caption{The closest systems against the four capabilities FL-IDS combines (\checkmark: a feature of the work; --: not a feature of the work)}
\label{tab:litcompare}
\end{table}

\section{Research Gaps}
From the survey, the following gaps were identified.
\begin{enumerate}
  \item \textbf{No single system combines all five parts.} Federated anomaly detection, cosine-based robust aggregation, a boosting-first cascade and FL with SDN enforcement each exist in some paper, as pairs or triples. A handful of targeted searches found no system combining all five. (This is not a systematic review, so novelty is stated cautiously.)
  \item \textbf{Zero-day detection is rarely measured per class.} Most federated IDS papers report accuracy on known classes only. The value of an anomaly-detection backstop shows up only when each attack class is held out in turn.
  \item \textbf{Robust aggregation under non-IID data.} Cosine-similarity defences assume honest updates are similar to each other. Under strong label skew, honest clients disagree too, and the literature reports that such defences degrade. This needs to be measured, not assumed.
  \item \textbf{Inflated results on the released Edge-IIoTset CSV.} Work that trains directly on the authors' preprocessed CSV inherits its formatting artefacts. The random packet-level split also puts packets of the same connection in both training and test sets.
  \item \textbf{Weak detection of application-layer attacks.} With only the released packet-header features, SQL injection, XSS, file upload, password attacks and HTTP floods are hard to tell apart from each other and from benign traffic, because most of their packets are content-free TCP acknowledgements.
  \item \textbf{Detection rarely drives verified enforcement.} Many papers stop at a classification score. Few show a detection producing a flow rule that is then read back from the switch.
\end{enumerate}

\section{Mapping of Project Objectives with Research Gaps}
\begin{table}[H]
\centering
\renewcommand{\arraystretch}{1.3}
\begin{tabular}{|p{6.2cm}|p{7.3cm}|}
\hline
\textbf{Research gap} & \textbf{Objective addressing it} \\ \hline
1. No system combines all five parts & All four objectives together: trust-filtered federated learning (1, 2), a compact model (3) and an autoencoder backstop (4), integrated with a boosting classifier and SDN mitigation \\ \hline
2. Zero-day detection rarely measured & Objective 4: autoencoder-based anomaly detection, measured by holding out each of the 14 attack classes \\ \hline
3. Robust aggregation under non-IID data & Objective 1: cosine/MAD/EMA-trust filtering of client updates, tested from 0\% to 50\% malicious clients \\ \hline
4. Inflated results on the released CSV & Supporting work for all objectives: placeholder-leak fix, MITM capture repair and a flow-grouped evaluation split \\ \hline
5. Weak application-layer detection & Supporting work: service-port, content and flow features for the classifier (Section~\ref{sec:newfeat}) \\ \hline
6. Detection not linked to enforcement & Objective 2 (practical FL design): Phase B turns a verdict into an OpenFlow rule verified in the switch's flow table \\ \hline
\end{tabular}
\caption{Mapping of objectives to research gaps}
\end{table}
